% foundations/wzw_group.tex
% Part of the Hopfion Framework Repository
% Source: Paper I — The Density-Feedback Faddeev-Niemi Hopfion: Exact Algebraic Predictions for Standard-Model Parameters
% Status: published
% Last updated: 2026-06-29

%% hopfion-framework/preamble.tex
%% Shared preamble for all repository files.
%% \input this at the top of any standalone document.
%% Repository files are designed to be \input into papers directly.

\usepackage{amsmath,amssymb,amsthm}
\usepackage{hyperref}
\usepackage{enumitem}
\usepackage{booktabs}
\usepackage{xcolor}
\usepackage{geometry}
\geometry{margin=2.5cm}

%── Theorem environments ──────────────────────────────────────────────────────
\newtheorem{theorem}{Theorem}[section]
\newtheorem{lemma}[theorem]{Lemma}
\newtheorem{proposition}[theorem]{Proposition}
\newtheorem{corollary}[theorem]{Corollary}
\theoremstyle{definition}
\newtheorem{definition}[theorem]{Definition}
\newtheorem{result}[theorem]{Result}
\theoremstyle{remark}
\newtheorem{remark}[theorem]{Remark}
\newtheorem{observation}[theorem]{Observation}
\newtheorem{conjecture}{Conjecture}
\newtheorem{construction}[theorem]{Construction}
\newtheorem{condition}[theorem]{Condition}

%── Repository metadata commands ─────────────────────────────────────────────
% Usage: \status{published} or \status{open} or \status{conjectured}
\newcommand{\status}[1]{\par\noindent{\color{blue}\textbf{Status:} #1}}
% Usage: \source{Paper I, Theorem thm:scale}
\newcommand{\source}[1]{\par\noindent{\color{darkgray}\textbf{Source:} #1}}
% Usage: \residual{0.013\%} or \residual{exact}
\newcommand{\residual}[1]{\par\noindent{\color{teal}\textbf{Residual:} #1}}
% Usage: \depends{foundations/pentagon, foundations/suppression}
\newcommand{\depends}[1]{\par\noindent{\color{purple}\textbf{Depends on:} #1}}

%── Framework macros ─────────────────────────────────────────────────────────
\newcommand{\vp}{\varphi}
\newcommand{\bst}{\beta^{*}}
\newcommand{\Seff}{S_{\mathrm{eff}}}
\newcommand{\TCMB}{T_{\mathrm{CMB}}}
\newcommand{\Lcond}{\Lambda_{\mathrm{cond}}}
\newcommand{\LUV}{\Lambda_{\mathrm{UV}}}
\newcommand{\LQCD}{\Lambda_{\mathrm{QCD}}}
\newcommand{\MPl}{M_{\mathrm{Pl}}}
\newcommand{\ERy}{E_{\mathrm{Ry}}}
\newcommand{\aem}{\alpha_{\mathrm{em}}}
\newcommand{\sinW}{\sin^{2}\!\theta_{W}}
\newcommand{\vEW}{v_{\mathrm{EW}}}
\newcommand{\twoI}{2I}
\newcommand{\twoT}{2T}
\newcommand{\twoO}{2O}
\newcommand{\Ehat}{\hat{E}_8}
\newcommand{\QH}{Q_H}
\newcommand{\Ztwo}{\mathbb{Z}_2}
\newcommand{\Zthree}{\mathbb{Z}_3}
\newcommand{\SU}{\mathrm{SU}}
\newcommand{\rCMB}{\rho_{\mathrm{CMB}}}
\newcommand{\Lobs}{\Lambda_{\mathrm{obs}}}
\newcommand{\rinfty}{\rho_{\infty}}
\newcommand{\oBD}{\omega_{\mathrm{BD}}}
\newcommand{\xNMC}{\xi_{\mathrm{NMC}}}
\newcommand{\Efb}{E_{\mathrm{fb}}}
\newcommand{\Kfb}{K_{\mathrm{fb}}}
\newcommand{\dimq}[1]{d_{#1}}
\newcommand{\deff}{d_{\mathrm{eff}}}
\newcommand{\Kth}{K_{\mathrm{th}}}
\newcommand{\Ja}{\mathcal{J}_a}
\newcommand{\Jfour}{\mathcal{J}_4}
\newcommand{\Jiso}{\mathcal{J}_{\mathrm{iso}}}
\newcommand{\rhoinf}{\rho_{\infty}}
\newcommand{\rhoa}{\rho_{\mathrm{atom}}}
\newcommand{\Heff}{H_{\mathrm{eff}}}
\newcommand{\epsr}{\varepsilon(r)}
\newcommand{\Hcond}{H_{\mathrm{cond}}}
\newcommand{\eps}{\varepsilon}
\newcommand{\Rzero}{R_0}
\newcommand{\Cstar}{C^*}
\newcommand{\mH}{m_H}
\newcommand{\aem}{\alpha_{\mathrm{em}}}
\newcommand{\ms}{\mu^*}
\newcommand{\lam}{\lambda}
\newcommand{\fzero}{f_0}
\newcommand{\ftwo}{f_2}
\newcommand{\Itor}[1]{I_{#1}^{\mathrm{tor}}}
\newcommand{\Jfb}{\mathcal{J}_{\mathrm{fb}}}
\newcommand{\bz}{\beta_0}
\newcommand{\sopt}{s_{\mathrm{opt}}}
\newcommand{\leff}{\lambda_{\mathrm{eff}}}
\newcommand{\ord}{\operatorname{ord}}
\newcommand{\G}{\mathcal{G}}
\newcommand{\Ia}{I_{2a}}
\newcommand{\Ib}{I_{4}}
\newcommand{\kk}{\kappa}
\newcommand{\aInv}{\alpha^{-1}}
\newcommand{\mWZW}{m_{\mathrm{WZW}}}
\newcommand{\tang}{\dot\Gamma}
\newcommand{\Jones}{V}
\newcommand{\lk}{\mathrm{lk}}
\newcommand{\phistar}{\phi^*}
\newcommand{\CC}{\mathbb{C}}
\newcommand{\HH}{\mathbb{H}}
\newcommand{\OO}{\mathbb{O}}
\newcommand{\ZZ}{\mathbb{Z}}
\newcommand{\pcond}{p_{\mathrm{cond}}}
\newcommand{\pQM}{p_{\mathrm{QM}}}
\newcommand{\xcond}{x_{\mathrm{cond}}}
\newcommand{\CHSH}{\mathrm{CHSH}}
\newcommand{\Tsirelson}{2\sqrt{2}}
\newcommand{\RR}{\mathbb{R}}
\newcommand{\Hil}{\mathcal{H}}
\newcommand{\CQ}{\mathcal{C}_Q}
\newcommand{\CQr}{\mathcal{C}_Q^{\mathrm{red}}}
\newcommand{\bOmega}{\boldsymbol{\Omega}}
\newcommand{\dmu}{d\mu}
\newcommand{\nlvl}[1]{n\!\left(#1\right)}
\newcommand{\IE}{\mathrm{IE}}
\newcommand{\bsF}{{}_2F_1}




\section*{WZW Group Theory and Exact Algebraic Chain}
\label{sec:wzw_group}

% Results extracted from Paper I
% Scan date: 2026-06-29

%════════════════════════════════════════════════════════════════════
% Group Theory — McKay Correspondence
%════════════════════════════════════════════════════════════════════

\begin{theorem}[McKay correspondence]
\label{P1:thm:mckay}
The McKay graph of $2I$ constructed from the two-dimensional irrep $\Gamma_2$
is the affine Dynkin diagram $\hat{E}_8$, establishing the chain
\begin{equation}
  2I \;\longleftrightarrow\; \hat{E}_8 \;\supset\;
  E_8 \;\supset\; \mathrm{SU}(5) \;\supset\;
  \mathrm{SU}(3)_C\times\mathrm{SU}(2)_L\times U(1)_Y.
  \label{P1:eq:mckay_chain}
\end{equation}
\end{theorem}
\status{published}
\source{P1:thm:mckay; McKay (1980)}

%════════════════════════════════════════════════════════════════════
% WZW Level and Quantum Dimension
%════════════════════════════════════════════════════════════════════

\begin{theorem}[WZW level $k=3$]
\label{P1:thm:k3}
The subgroup $\mathbb{Z}_{10}\subset 2I$ forces the $\mathrm{SU}(2)$ WZW
model to level $k=3$ uniquely via
\begin{equation}
  \ord(q) = 2(k+2)\big|_{k=3} = 10 = |\mathbb{Z}_{10}|,
  \qquad
  q = e^{i\pi/(k+2)} = e^{i\pi/5}.
  \label{P1:eq:q}
\end{equation}
\end{theorem}
\status{published}
\source{P1:thm:k3}

\begin{theorem}[Quantum dimension equals $\vp$]
\label{P1:thm:dimq}
The quantum dimension of the spin-$\tfrac{1}{2}$ primary at WZW level $k=3$
equals $\vp$:
\begin{equation}
  \dim_q\!\left(\tfrac{1}{2};\,k=3\right)
  = \frac{S_{0,\,1/2}}{S_{0,0}}
  = \frac{\sin(2\pi/5)}{\sin(\pi/5)}
  = 2\cos(\pi/5)
  = \vp,
  \label{P1:eq:dimq}
\end{equation}
where $S$ is the modular $S$-matrix of $\mathrm{SU}(2)_3$.  This
coincides with the character value $\chi_{\Gamma_2}(C_5)=\vp$ of the
$2I$ representation and is the unique fixed point of $x^2 = x+1$.
\end{theorem}
\status{published}
\source{P1:thm:dimq}
\residual{exact}

%════════════════════════════════════════════════════════════════════
% Exact Algebraic Chain
%════════════════════════════════════════════════════════════════════

\begin{theorem}[Scale selection]
\label{P1:thm:scale}
At the self-consistent $\bst$,
\begin{equation}
  \sopt = 2^{1/6}
  \quad\Longleftrightarrow\quad
  \frac{\Jfour}{\Ja} = \frac{2^{4/3}}{\vp^{5}}.
  \label{P1:eq:equiv}
\end{equation}
These two conditions are algebraically identical; neither is more primitive
than the other.
\end{theorem}
\status{published}
\source{P1:thm:scale}
\begin{proof}
From $\sopt^{2} = \vp^{6}\Jfour/(2\vp\,\Ja) = (\vp^{5}/2)\cdot(\Jfour/\Ja)$
(using $\Kfb=2\vp\Ja$ and $\lam=\vp^6$):
$\sopt=2^{1/6}\Leftrightarrow\sopt^{6}=2\Leftrightarrow
(\vp^{5}/2)^{3}(\Jfour/\Ja)^{3}=2\Leftrightarrow\Jfour/\Ja = 2^{4/3}/\vp^{5}$.
\end{proof}
\status{published}
\source{P1:thm:scale}
\residual{exact (algebraic identity); numerically $0.020\%$ at $h=0.12$}
\depends{profile/virial.tex (prop:Kfb); electroweak/weinberg.tex (P1:cor:lambda)}

\begin{theorem}[Charge quantisation]
\label{P1:thm:charge_q10}
The topological charge of the Hopfion condensate at the CMB scale is
\begin{equation}
  Q \;=\; \ord(q_{2I}) \;=\; 2(k+2)\big|_{k=3} \;=\; 10 \quad\text{(exact)}
  \label{P1:eq:Q_q10}
\end{equation}
and is realised numerically as
\begin{equation}
  Q \;=\; \bz\cdot\rCMB \;=\; \vp^{9}\sqrt{\Ja\cdot\Jfour} \;=\; 10.012
  \quad(0.12\%),
  \label{P1:eq:br10}
\end{equation}
where $\bz = \TCMB(\pi^2/150)^{1/4} \approx 119\;\mu\mathrm{eV}$ and
$\rCMB = (\pi^2/15)\TCMB^4$ is the CMB energy density.
\end{theorem}
\status{published}
\source{P1:thm:charge_q10}
\residual{exact (group-theoretic route); $0.12\%$ (numerical route, lattice artefact closed in Paper~XII)}
\depends{foundations/wzw_group.tex (P1:thm:k3)}

%════════════════════════════════════════════════════════════════════
% Physical Interpretations
%════════════════════════════════════════════════════════════════════

\begin{remark}[Physical interpretation of the feedback kernel]
\label{P1:rem:chameleon}
The functional $\Jfb = \int \mathrm{kern}/(1+\bst\,\mathrm{kern})\,\mathrm{dvol}$
admits three complementary physical interpretations.

\textbf{(i) Born--Infeld regularisation.}
$\Jfb$ is the Born--Infeld saturation of the isotropic FN kinetic
energy $\Jiso$: as $\mathrm{kern}\to\infty$ the integrand saturates at
$1/\bst$, preventing ultraviolet divergence of the gradient energy.
This interpretation is a derivation via \texttt{thm:vacuum}: $\Jfb$ is
exactly the gradient energy of the parent Lagrangian $\mathcal{L}_{\mathrm{kin}}
= |\nabla\rho|^2/[\vp^6(1+\bst|\nabla\rho|^2)]$
evaluated on the $Q=2$ Hopfion sector.

\textbf{(ii) Chameleon scalar field.}
A scalar field coupled to ambient matter acquires a density-dependent
effective mass $m_{\mathrm{eff}}^2 \propto \rho_{\mathrm{matter}}$,
screening short-wavelength fluctuations in dense regions.
Here $\mathrm{kern}$ plays the role of local field energy density,
and $1/(1+\bst\,\mathrm{kern})$ is the screening factor.
The self-consistency condition $\Jfb/\Ja = \vp$ is the field-theoretic
fixed-point equation analogous to the chameleon equilibrium condition.

\textbf{(iii) Madelung quantum-pressure correction.}
In the weak-field expansion,
$\mathrm{kern}/(1+\bst\,\mathrm{kern}) = \mathrm{kern} - \bst\,\mathrm{kern}^2 + O(\bst^2)$,
the correction term $-\bst\,\mathrm{kern}^2$ has the same structure as
the quantum-pressure ($D^2k^4$) correction in the Madelung--Bogoliubov
dispersion relation.
\end{remark}
\status{published}
\source{P1:rem:chameleon}

\begin{remark}[Topological stability of the feedback functional]
\label{P1:rem:stability}
Like the standard FN energy, $\Efb = \Kfb\cdot\Jfour$ has no
Derrick scale-barrier: both $\Kfb$ and $\Jfour$ decrease as $f\to 0$,
so the trivial vacuum $f=0$ is the global minimum of $\Efb$.
The Hopfion is the saddle-point minimum only within the $Q=2$ topological
sector.  The saddle-snapshot method (Definition~\ref{P1:def:snapshot})
circumvents this instability by optimising the scale-invariant proxy
$\Efb$ and recording the step at which $\sopt=1$.
\end{remark}
\status{published}
\source{P1:rem:stability}

\begin{remark}[Connection to $E_8$ and WZW]
\label{P1:rem:e8}
The value $\sopt = 2^{1/(2k)}$ with $k=3$ follows from the Coxeter number
of $E_8$: $h(E_8) = 30 = 2k(k+2)$ is satisfied uniquely by $k=3$ among
ADE algebras with $h\leq 30$.  The WZW level $k=3$ is forced by
$\mathbb{Z}_{10}\subset 2I$ (Theorem~\ref{P1:thm:k3}).  The factor of~$2$
in $\sopt^{6}=2$ is $|2I|/|I|=2$, the order of the central extension.
\end{remark}
\status{published}
\source{P1:rem:e8}

%════════════════════════════════════════════════════════════════════
% Vacuum Uniqueness
%════════════════════════════════════════════════════════════════════

\begin{theorem}[Vacuum uniqueness at $\Lcond$]
\label{P1:thm:vacuum}
Among all static finite-action configurations of the density-feedback
FN functional at the condensate scale $\Lcond = \TCMB(\pi^2/150)^{1/4}$,
the $Q=2$ icosahedral Hopfion is the unique minimum-action configuration
within its topological sector, and has strictly lower action than all
$Q=0$ and $Q=1$ competitors.  Specifically:
\begin{enumerate}[noitemsep,label=\emph{(\roman*)}]
  \item \textit{$Q=0$ (uniform):} Any non-trivial gradient
    knot with gradient concentrated near $\theta=0$ has strictly lower
    angular cost than the isotropic configuration.
  \item \textit{$Q=0$ (spherically symmetric):} A spherically symmetric
    profile saturates the angular cost; any $Q=2$ Hopfion with toroidal
    geometry achieves the same topological energy $\Jfour$ with lower
    angular cost.
  \item \textit{$Q=1$ (single-tube):} Two single-tube quanta fuse via
    $j_{1/2}\times j_{1/2} = j_0\oplus j_1$ in $\mathrm{SU}(2)_3$.
    The singlet channel $j_0$ (the $Q=2$ condensate) has conformal
    weight $h_0=0$ strictly less than the triplet $h_1 = 2/5$, so the
    $Q=2$ state is lower energy.
  \item \textit{$Q\geq 3$:} The Vakulenko--Kapitanski bound
    $E\geq C_{\mathrm{VK}}|Q|^{3/4}$ gives a supralinear energy cost.
  \item \textit{Icosahedral geometry among $Q=2$:} The binary icosahedral
    group $2I\subset\mathrm{SU}(2)$ has order $|2I|=120$, the maximum
    among binary polyhedral groups.  The condensate temperature uniquely
    selects $2I$ via $\bz\rCMB = \ord(q_{2I})=10$
    (Theorem~\ref{P1:thm:charge_q10}).
\end{enumerate}
\end{theorem}
\status{published}
\source{P1:thm:vacuum}
\begin{proof}
Parts (i)--(ii) follow from the inequality
$\int S_{\mathrm{eff}}^{\mathrm{knot}}\,d\Omega <
\int S_{\mathrm{eff}}^{\mathrm{iso}}\,d\Omega$
for configurations whose gradients concentrate near $\theta=0$.
Part (iii) is the conformal weight comparison $h_0 = 0 < h_1 = 2/5$
in $\mathrm{SU}(2)_3$.
Part (iv) follows directly from the Vakulenko--Kapitanski inequality.
Part (v) is Theorem~\ref{P1:thm:charge_q10} together with the maximality of
$|2I|$ among binary polyhedral groups.
\end{proof}
\status{published}
\source{P1:thm:vacuum}
\depends{foundations/wzw_group.tex (P1:thm:charge_q10, P1:thm:k3)}

\begin{remark}
\label{P1:rem:symmetry_breaking_noncircular}
The symmetry breaking in part~(i) is not circular.
$S_{\mathrm{eff}}(\theta,\rho)$ is introduced as the suppression
function of the parent Lagrangian~\eqref{P1:eq:Lkin}, defined
\emph{before} any reference to a specific vacuum.  It breaks the
$O(3)$ symmetry of the free FN functional to $U(1)$ by its $\theta$-dependence,
and the proof shows this $U(1)$ remnant favours the toroidal Hopfion
over all $O(3)$-symmetric competitors.  The conclusion is that it is, within the
density-feedback framework, a consequence of $S_{\mathrm{eff}}$
and the condensate temperature.
\end{remark}
\status{published}
\source{P1:rem:symmetry_breaking_noncircular}

%════════════════════════════════════════════════════════════════════
% Witten Bosonization and Verlinde Formula — Paper III
%════════════════════════════════════════════════════════════════════

\begin{theorem}[Witten bosonization dictionary, $k=3$]
\label{P3:thm:bosonization}
The following bosonic and fermionic descriptions are exactly equivalent:
\begin{center}
\begin{tabular}{ll}
\toprule
Bosonic (WZW / Hopfion) & Fermionic ($k=3$ Dirac) \\
\midrule
$\mathrm{SU}(2)_3$ WZW field $g(x)$ & 3 free Dirac fermions $\psi_a$ \\
Curvature $\kk_{\mathrm{iso}}$ & Fermion density $\rho_F$ \\
$J_{2a} = \int\!\sin^4\!f\cdot\kk\,dV$ & Topological charge $Q_F$ \\
$\Jfb(\beta) = \int\!\kk/(1+\beta\kk)\,dV$ & Resolvent density $G(\beta)$ \\
WZ coupling $\mu^* = \sqrt{k+2}/\vp$ & WZ coupling of $k=3$ system \\
\bottomrule
\end{tabular}
\end{center}
\end{theorem}
\status{published}
\source{P3:thm:bosonization}
\depends{foundations/wzw_group.tex (P1:thm:k3, P1:thm:dimq); foundations/suppression.tex (P3:lem:bridge)}

\begin{remark}[The 3D/2D reduction]
\label{P3:rem:3d2d_reduction}
Witten's theorem applies to the 2D cross-section of the Hopfion torus.
In the thin-torus limit ($R_0\to\infty$), the 3D Hopfion problem reduces
exactly to a 2D problem on the tube cross-section, and the bosonization
is exact.  For finite $R_0$, corrections are $O(R_0^{-2})$, matching
the rate established in Proposition~\ref{P3:prop:virial_limit}.
This remark is promoted to a theorem with explicit error bounds below
(Theorem~\ref{P3:thm:bos_3d}).
\end{remark}
\status{published}
\source{P3:rem:3d2d_reduction}

\begin{theorem}[3+1D Witten bosonization with error bound]
\label{P3:thm:bos_3d}
At energies $E\ll r_{\mathrm{core}}^{-1}$, the 3+1D condensate $k$-essence
action is equivalent to the $\mathrm{SU}(2)_3$ WZW action plus three free
Dirac fermions, with corrections bounded by
$\|S_{3+1\mathrm{D}} - S_{\mathrm{WZW}}\| \leq C(Er_{\mathrm{core}})^2 R_0^{-2}$.
The WZ term descends from the 3+1D topological term via Zumino descent, and
the $\mathrm{SU}(2)_3$ level $k=3$ is preserved exactly.
\end{theorem}
\status{proved (via KK reduction, Zumino descent, anomaly matching)}
\source{P3:thm:bos_3d}

\begin{corollary}[Physical validity of the bosonization dictionary]
\label{P3:cor:bos_valid}
At the condensate linking scale, $\|S_{3+1\mathrm{D}}-S_{\mathrm{WZW}}\|
\leq C\,Q^{-2/3}R_0^{-2}$.  For $Q=10$ in the thin-torus limit
$R_0\gg 1$: the correction is $O(R_0^{-2})\ll 1$.
All SM-sector predictions of Theorem~\ref{P3:thm:bosonization} hold as exact
results in the limit $R_0\to\infty$, with fractional corrections $O(R_0^{-2})$.
\end{corollary}
\status{published}
\source{P3:cor:bos_valid}

\begin{remark}[Scope of the bosonization theorem]
\label{P3:rem:bos3d_scope}
Theorem~\ref{P3:thm:bos_3d} proves: (i) the 3+1D action reduces to 2D WZW at
low energies with explicit $O(R_0^{-2})$ bounds; (ii) the WZ term descends
correctly with level $k=3$ preserved; (iii) Hopf charge $Q=10$ is preserved
under KK reduction; (iv) 't~Hooft anomaly matching holds.
Not addressed: finite-$R_0$ corrections beyond $O(R_0^{-2})$, and
convergence of the KK expansion for non-static configurations.
\end{remark}
\status{published}
\source{P3:rem:bos3d_scope}

\begin{proposition}[$Q=2$ sector assignment]
\label{P3:prop:q2_sector}
The $Q=2$ Hopfion is the singlet ($j=0$) bound state of two $Q=1$ Hopfions
in the fusion decomposition $j_{1/2}\otimes j_{1/2} = j_0\oplus j_1$ of
$\mathrm{SU}(2)_3$.  The singlet has conformal dimension $h_0=0 < h_1=2/5$
and is the physical ground state.
\end{proposition}
\status{published}
\source{P3:prop:q2_sector}

\begin{theorem}[Verlinde/Cardy formula gives $V^*=\vp$]
\label{P3:thm:verlinde}
In the $j=0$ sector of $\mathrm{SU}(2)_3$, the vacuum expectation value of the
primary operator $\phi_{1/2}$ is
\begin{equation}
  \langle\phi_{1/2}\rangle_{j=0}
  = \frac{S_{0,1/2}}{S_{0,0}}
  = \frac{\sin(2\pi/5)}{\sin(\pi/5)}
  = 2\cos(\pi/5) = \vp.
\end{equation}
\end{theorem}
\status{published}
\source{P3:thm:verlinde}
\begin{proof}
The Verlinde $S$-matrix for $\mathrm{SU}(2)_3$ at $k+2=5$ gives
$S_{0,1/2}/S_{0,0} = \sin(2\pi/5)/\sin(\pi/5) = 2\cos(\pi/5) = \vp$
(pentagon identity).
\end{proof}
\status{published}
\source{P3:thm:verlinde}
\residual{exact}

\begin{theorem}[Bosonization proof of $V^*=\vp$]
\label{P3:thm:bos_main}
Under the Witten bosonization (Theorem~\ref{P3:thm:bosonization}), the
density-feedback virial satisfies $V^* = \Jfb(\bst)/\Ja = \vp$,
where $\bst$ is the self-consistent feedback parameter.
\end{theorem}
\status{published}
\source{P3:thm:bos_main}
\begin{proof}[Proof sketch]
Under the WZW/fermion duality, $V = \Jfb/\Ja \leftrightarrow \langle\phi_{1/2}(\beta)\rangle$.
At the self-consistent $\bst$ (the WZW fixed point), Theorem~\ref{P3:thm:verlinde}
gives $V^* = S_{0,1/2}/S_{0,0} = \vp$.
\end{proof}
\status{published}
\source{P3:thm:bos_main}
\residual{exact}
\depends{foundations/wzw_group.tex (P3:thm:bosonization, P3:thm:verlinde)}
